\documentclass{beamer}

\usepackage[utf8]{inputenc}
\usepackage[brazil]{babel}
\usepackage{graphicx}

\usetheme{Madrid}

\title{Análise de Programação Paralela\\
na Busca de Substrings}
\author{Samuel Minuzzo Pereira \\ Edilson Daniel de Souza}
\institute{UTFPR}
\date{\today}

\begin{document}



\frame{\titlepage}


\begin{frame}{Introdução}

\textbf{Busca de padrões em texto} é um problema clássico presente em:

\begin{itemize}
    \item indexação de documentos e sistemas de busca
    \item análise de logs e monitoramento de sistemas
    \item bioinformática e análise forense
\end{itemize}

\vspace{0.4cm}

\textbf{Motivação prática:}
\begin{itemize}
    \item arquivos de texto podem chegar a \textbf{múltiplos gigabytes}
    \item abordagens simples tornam-se inviáveis nessa escala
    \item necessidade de algoritmos eficientes \textbf{e} paralelização
\end{itemize}

\vspace{0.4cm}

\textbf{Este trabalho} avalia, na prática, o impacto de diferentes
escolhas de implementação no desempenho computacional.

\end{frame}


\begin{frame}{Objetivo}

\textbf{Objetivo geral:} implementar, executar e comparar quatro versões
do problema de busca de substrings em arquivos de grande porte.

\vspace{0.4cm}

\begin{columns}
    \column{0.5\textwidth}
    \textbf{Versões implementadas:}
    \begin{itemize}
        \item Java sequencial
        \item Java paralelo (threads)
        \item C sequencial
        \item C paralelo (OpenMP)
    \end{itemize}

    \column{0.5\textwidth}
    \textbf{Métricas avaliadas:}
    \begin{itemize}
        \item tempo de execução
        \item speedup
        \item corretude dos resultados
    \end{itemize}
\end{columns}

\vspace{0.5cm}

\textbf{Questões centrais:}
\begin{itemize}
    \item Quanto a paralelização melhora o desempenho?
    \item Qual linguagem se sai melhor e por quê?
\end{itemize}

\end{frame}


\begin{frame}{Problema}

\textbf{Definição:} contar todas as ocorrências de uma substring $P$
em um arquivo texto $T$, incluindo ocorrências sobrepostas.

\vspace{0.3cm}

\textbf{Formalmente:} determinar o número de índices $i$ tais que:
\[
T[i\,..\,i+m-1] = P, \quad 0 \leq i \leq n-m
\]

onde $n = |T|$ e $m = |P|$.

\vspace{0.3cm}

\begin{columns}
    \column{0.5\textwidth}
    \textbf{Entrada:}
    \begin{itemize}
        \item caminho do arquivo texto
        \item substring buscada
        \item número de threads
    \end{itemize}

    \column{0.5\textwidth}
    \textbf{Saída:}
    \begin{itemize}
        \item total de ocorrências
        \item tempo sequencial
        \item tempo paralelo
        \item speedup
    \end{itemize}
\end{columns}

\vspace{0.3cm}

\small\textit{O texto é tratado como sequência de bytes,
tornando a solução independente de codificação.}

\end{frame}


\begin{frame}{Complexidade}

\begin{columns}
    \column{0.5\textwidth}
    \textbf{Busca ingênua:}
    \begin{itemize}
        \item alinha $P$ em cada posição de $T$
        \item compara caractere a caractere
        \item complexidade: $O(n \cdot m)$
        \item inviável para arquivos de GBs
    \end{itemize}

    \column{0.5\textwidth}
    \textbf{Boyer-Moore-Horspool:}
    \begin{itemize}
        \item pré-processa o padrão $P$
        \item permite \textbf{saltar posições} do texto
        \item pior caso ainda $O(n \cdot m)$
        \item mas muito mais eficiente na prática
    \end{itemize}
\end{columns}

\vspace{0.5cm}

\textbf{Por que BMH e não KMP ou Boyer-Moore clássico?}
\begin{itemize}
    \item boa relação entre \textbf{simplicidade} e \textbf{desempenho prático}
    \item adaptação natural para leitura em blocos (arquivos grandes)
    \item tratamento direto byte a byte
\end{itemize}

\end{frame}


\begin{frame}{Boyer-Moore-Horspool}

\textbf{Pré-processamento:} tabela de deslocamento para cada byte possível
\[
shift[P[i]] = m - 1 - i, \quad 0 \leq i < m-1
\]
Posições não presentes no padrão recebem deslocamento $m$.

\vspace{0.2cm}

\textbf{Funcionamento a cada etapa:}
\begin{enumerate}
    \item alinha $P$ a uma posição do texto
    \item compara do \textbf{último} caractere para o primeiro
    \item se coincidir totalmente: ocorrência encontrada
    \item se não: consulta a tabela e \textbf{salta} posições
\end{enumerate}

\vspace{0.2cm}

\end{frame}

\begin{frame}
\textbf{Salto pós-match via \textit{KMP failure function}:}

\begin{columns}
    \column{0.5\textwidth}
    \textbf{Padrão sem sobreposição:}\\
    \texttt{Bacamarte} ($m = 9$)
    \begin{itemize}
        \item nenhum prefixo é sufixo
        \item salto = $m = 9$
        \item máximo aproveitamento
    \end{itemize}

    \column{0.5\textwidth}
    \textbf{Padrão com sobreposição:}\\
    \texttt{aa} ($m = 2$) em \texttt{aaaa}
    \begin{itemize}
        \item prefixo \texttt{a} também é sufixo
        \item salto = $1$ (não perde matches)
        \item ocorrências: 3
    \end{itemize}
\end{columns}



\end{frame}


\begin{frame}{Paralelização}

\textbf{Desafio central:} o texto não pode ser dividido em partes
independentes sem cuidado nas fronteiras.

\vspace{0.2cm}

\textbf{Problema das fronteiras:}
\begin{itemize}
    \item uma ocorrência pode \textbf{iniciar} num bloco e \textbf{terminar} no próximo
    \item divisão ingênua causa perda ou dupla contagem de ocorrências
\end{itemize}

\vspace{0.2cm}

\end{frame}

\begin{frame}


\textbf{Soluções adotadas:}
\begin{columns}
    \column{0.5\textwidth}
    \textbf{Sobreposição de blocos:}
    \begin{itemize}
        \item cada bloco lê $m-1$ bytes extras do bloco seguinte
        \item garante que nenhuma ocorrência seja perdida
    \end{itemize}

    \column{0.5\textwidth}
    \textbf{Intervalo de propriedade:}
    \begin{itemize}
        \item cada thread inspeciona além do seu limite
        \item mas contabiliza \textbf{apenas} ocorrências cujo índice inicial pertence ao seu intervalo
    \end{itemize}
\end{columns}

\vspace{0.2cm}

\textbf{Double buffering (C paralelo):}
\begin{itemize}
    \item enquanto uma thread processa o bloco atual, outra já lê o próximo do disco
    
\end{itemize}


\end{frame}

\begin{frame}{Implementações}

\textbf{Núcleo comum a todas as versões:}
\begin{itemize}
    \item algoritmo Boyer-Moore-Horspool
    \item leitura em blocos (evita carregar GBs em memória)
    \item sobreposição de $m-1$ bytes entre blocos consecutivos
    \item salto pós-match via \textit{KMP failure function}
\end{itemize}

\vspace{0.3cm}

\begin{columns}
    \column{0.5\textwidth}
    \textbf{Java:}
    \begin{itemize}
        \item leitura: \texttt{BufferedInputStream}
        \item paralelismo: \texttt{ExecutorService}
        \item overhead da JVM e GC
        \item maior portabilidade
    \end{itemize}

    \column{0.5\textwidth}
    \textbf{C:}
    \begin{itemize}
        \item leitura binária com buffer manual
        \item paralelismo: \texttt{OpenMP}
        \item controle direto de memória
        \item double buffering (versão paralela)
    \end{itemize}
\end{columns}

\vspace{0.3cm}

\end{frame}

\begin{frame}{Ambiente de Testes}

\begin{columns}
    \column{0.5\textwidth}
    \textbf{Hardware:}
    \begin{itemize}
        \item \textbf{Modelo:} Acer Predator Helios Neo 16
        \item \textbf{CPU:} Intel Core i7-13650HX
        \item \textbf{Núcleos físicos:} 14
        \item \textbf{Threads lógicas:} 20
        \item \textbf{RAM:} 24 GB
        \item \textbf{Disco:} SSD NVMe Kingston\\1 TB
    \end{itemize}

    \column{0.5\textwidth}
    \textbf{Software:}
    \begin{itemize}
        \item \textbf{SO:} Windows 11 Home 64 bits
        \item \textbf{Compilador C:} GCC 15.2.0
        \item \textbf{JVM:} OpenJDK 17.0.18 (Temurin)
        \item \textbf{Automação:} script Python
    \end{itemize}

    \vspace{0.4cm}

    \textbf{Threads testadas:}\\
    \vspace{0.1cm}
    1, 2, 8, 20, 40, 80, 120, 160, 200
\end{columns}

\end{frame}
\begin{frame}{Entradas}

\textbf{Tamanhos de arquivo testados:}\\
\vspace{0.1cm}
0,1 GB \quad 1 GB \quad 2 GB \quad 4 GB \quad 8 GB \quad 16 GB

\vspace{0.4cm}

\textbf{Substrings escolhidas e sua relação com o BMH:}

\vspace{0.2cm}

\begin{tabular}{cll}
\hline
\textbf{Substring} & \textbf{Característica} & \textbf{Comportamento no BMH} \\
\hline
\texttt{a} & comprimento 1 & pior caso \\
\texttt{como} & comprimento 4 & caso típico\\
\texttt{Bacamarte} & comprimento 9 & melhor caso\\
\hline
\end{tabular}

\vspace{0.3cm}


\end{frame}



\begin{frame}{Resultados — Tempo de Execução (segundos)}

\centering
\footnotesize
\begin{tabular}{ccccc}
\hline
\textbf{Arquivo} & \textbf{Java Seq} & \textbf{Java Par} & \textbf{C Seq} & \textbf{C Par} \\
\hline
0,1 GB  & 0,159  & 0,121  & 0,078 & 0,082 \\
1 GB    & 1,383  & 0,765  & 0,686 & 0,633 \\
2 GB    & 2,483  & 1,324  & 1,263 & 1,051 \\
4 GB    & 4,875  & 2,617  & 2,512 & 1,934 \\
8 GB    & 9,530  & 4,932  & 4,878 & 3,663 \\
16 GB   & 34,980 & 32,979 & 28,087 & 31,699 \\
\hline
\end{tabular}

\vspace{0.4cm}

\end{frame}

\begin{frame}{Análise dos Resultados}

\begin{columns}
    \column{0.5\textwidth}
    \textbf{C vs. Java:}
    \begin{itemize}
        \item C mantém menores tempos absolutos em vários cenários
        \item Java apresentou melhor ganho relativo de paralelismo
        \item em 16 GB, Java ainda ganha pouco com paralelismo
        \item em 16 GB, C paralelo ficou pior que C sequencial
    \end{itemize}

    \vspace{0.3cm}

    \textbf{Escalabilidade com threads:}
    \begin{itemize}
        \item escalabilidade não linear
        \item melhores speedups surgem em faixas altas (80 a 160)
        \item ganho marginal cai com gargalos de I/O e sincronização
    \end{itemize}

    \column{0.5\textwidth}
    \textbf{Por que 16 GB quebrou o speedup?}
    \begin{itemize}
        \item arquivo excede cache efetivo do SO
        \item leituras diretas do SSD a cada chunk
        \item largura de banda do disco torna-se\\o gargalo dominante
        \item mais threads não reduzem tempo de I/O
    \end{itemize}

    \vspace{0.3cm}

    \textbf{Influência da substring:}
    \begin{itemize}
        \item \texttt{a}: pior desempenho (salto = 1)
        \item \texttt{como}: desempenho moderado
        \item \texttt{Bacamarte}: melhor desempenho\\(saltos grandes, padrão raro)
    \end{itemize}
\end{columns}

\end{frame}
\begin{frame}{Dificuldades Encontradas}

\begin{columns}
    \column{0.5\textwidth}

    \textbf{Fronteiras entre blocos:}
    \begin{itemize}
        \item ocorrência pode cruzar dois chunks
        \item solução: sobreposição de $m-1$ bytes
    \end{itemize}

    \vspace{0.3cm}

    \textbf{Dupla contagem entre threads:}
    \begin{itemize}
        \item threads adjacentes inspecionam\\regiões sobrepostas
        \item solução: intervalo de propriedade\\(cada thread conta apenas\\ocorrências do seu intervalo)
    \end{itemize}

    \column{0.5\textwidth}

    \textbf{Salto pós-match:}
    \begin{itemize}
        \item avanço 1: correto, mas lento
        \item avanço $m$: rápido, mas perde\\sobreposições
        \item solução: \textit{KMP failure function}\\calcula o salto mínimo seguro
    \end{itemize}

    \vspace{0.3cm}

    \textbf{Automação dos testes:}
    \begin{itemize}
        \item padronização de saída C e Java
        \item diferenças de \textit{locale} no Java\\(\texttt{Locale.US} na formatação)
        \item integração com script Python
    \end{itemize}

\end{columns}

\end{frame}

\begin{frame}{Conclusão}

\textbf{O que foi alcançado:}
\begin{itemize}
    \item implementação e comparação de 4 versões do problema
    \item algoritmo BMH com salto pós-match via \textit{KMP failure function}
    \item double buffering na versão C paralela
\end{itemize}

\vspace{0.3cm}

\textbf{Principais conclusões:}
\begin{itemize}
    \item C teve melhor tempo absoluto em boa parte dos cenários
    \item Java paralelo teve melhor ganho relativo (speedup médio)
    \item paralelismo é eficaz até 8 GB; em 16 GB o limite de I/O domina
    \item BMH adequado para o problema, especialmente com padrões longos
    \item não há uma solução única melhor para todos os critérios
\end{itemize}

\vspace{0.3cm}


\end{frame}


\begin{frame}
\centering
\Huge Obrigado!
\end{frame}


\end{document}